\section{Uniformidad de las constantes y rango de uso}\label{sec:usos}
\paragraph{Citas.} Además de las anteriores:
\begingroup\sloppy
\begin{itemize}
\item S02, S03, E01, M01, R05 y D04 son \path{S/02_menor.md}, \path{S/03_determinante.md},
\path{E/01_enunciado.md}, \path{M/01_enunciado.md}, \path{R/05_bonferroni.md} y
\path{D/04_difusion.md}, en \path{r19/expediente/originales/};
\item 04m es \path{r21/consolidado/latex/04_marginales.tex}, idéntico al de v2;
\item E02 es \path{fuentes/velocidad_via2_ronda8_20261116T072748Z/latex/02_expansion.tex};
\item R1m es \path{fuentes/velocidad_via2_ronda1_20261116T023056Z/latex/mascara.tex};
\item K06 y K08 son \path{06_torres.tex} y \path{08_estabilidad.tex}, en
\path{fuentes/velocidad_via2_ronda9_20261116T082611Z/latex/};
\item MM07 es \path{fuentes/velocidad_via2_media_matriz_20261116T150154Z/latex/07_vertices.tex};
\item GI02 es \path{fuentes/velocidad_via2_grupos_interlinea_20261116T161315Z/latex/02_exterior.tex};
\item W2d es \path{W2_difusion/ACTA.md}, y CONST es \path{CONSTANTES.md} de este consolidado,
cuyo paso~1 fija ya $C_*$.
\end{itemize}
\endgroup

\paragraph{1. Uniformidad en la clase.} $C_*$ sólo depende de $(A,r)$: a través de $(K_U,C_U,c_U)$,
que dependen de $(1+A,r)$ (U01:19--20, 43; E1:36), y de $\rho,\theta,c_H,L,C_1,C_2,C_S,C_E,C_x$,
explícitas en $(A,r)$. No depende de $K$, del soporte (se admiten ceros) ni de otra propiedad de
$D$, y vale para todo $K\ge2$: al aplicarlo no aparece ningún umbral de tamaño. Si $A\le A'$ y
$r\le r'$, entonces $\mathcal D_K(A,r)\subset\mathcal D_K(A',r')$, y basta acotar los márgenes.
Es una constante de clase: se fija con los márgenes (CONST:3--4), antes que $b$, $\rho$,
$\alpha$ y $n$, sin crear ningún círculo. Todo es determinista en una máscara fija; no interviene
ninguna ley de las historias.

\paragraph{2. Los usos.} En cada uso el Teorema se aplica a matrices reales, no negativas y
bistocásticas de una clase $\mathcal D_K(A,r)$ cuyos parámetros son de clase. Como
$|B_{ij}|\le\max(D_{ij},1/K)$, la cota $KD_{ij}\le M$ da $A=\max(M,1)$. En todos los usos, además,
$D$ es el escalado exacto de una máscara binaria, la de su soporte, como pide literalmente U01:9:
el escalado de Sinkhorn no cambia al multiplicar la entrada por diagonales positivas, y el menor
reescalado de un escalado es el escalado del menor de la máscara. Así, aun leyendo U01 en sentido
estricto, todos los usos quedan cubiertos. La pertenencia a la clase en cada uso la certifican E2,
W2\_torres, W2\_difusion, W2\_pares y W2\_t4, y la de U02, E1 (\path{W4_rev/ACTA.md}:263, 7b).
\begin{enumerate}
\item \emph{S02--S03 de E2} (04m:6--18). S03:10--15 lo aplica a $D$ y a $D'$.
  \begin{itemize}
  \item $D$ es el escalado de $P\circ(1-P)$ (S02:9--16). Como $P_{ij}(1-P_{ij})$ es producto de una
  función de fila y otra de columna en el soporte, es el escalado exacto $D_A$ de la máscara. Para
  $j\ge j_D$ cumple $jD\le4C$ y $\|D-J\|_{\rm op}\le1-\gamma_T/2$ (W2d:57, 84, 212); por tanto
  $D\in\mathcal D_j(\max(4C,1),1-\gamma_T/2)$.
  \item $D'$ es el menor reescalado de orden $j-2$: $\max D'=O(1/j)$ y $\|D'-J_{j-2}\|\le r'+O(1/j)$
  (S02:43--47; U02:35, 44--45, certificada por E1). Para $j\ge j_1$ está en
  $\mathcal D_{j-2}(A',(1+r')/2)$, con $A'$ de clase.
  \item En el telescopio de 04m:58--64 se aplica a los menores que prescriben
  $d\le A_0\log j/\log\log j$ sitios y a sus menores de dos sitios. Sus márgenes cambian por factores
  $1+O(d/j)$ y la brecha pierde $O(d/j)$ (R05:13--23; 04m:62). Para $j\ge j_2$ están en una clase
  común, con $K\in[j-d-2,\,j-d]$: el menor que prescribe $d$ sitios tiene tamaño $j-d$, y su menor
  de dos sitios, $j-d-2$ (\path{W4_rev/ACTA.md}:261, 7a).
  \end{itemize}
\item \emph{E01} (E02:14--23; R1m:2--8). Se aplica al escalado $D$ de $P\circ(1-P)$ de una máscara
  con $p\le M/k$ y $\|P-J\|_{\rm op}\le r$ (E01:9--12). Como $D=P(1+O(1/k))$ en el soporte y
  $\|D-P\|_{\rm HS}=O(1/k)$ (R1m:5--7), para $k\ge k_0(M,r)$ se tiene
  $D\in\mathcal D_k(\max(2M,1),(1+r)/2)$.
\item \emph{K08}.
  \begin{itemize}
  \item E01 sobre las máscaras replicadas de orden $bk$, con $b\to\infty$ (K08:2--9). Su escalado es
  $D^{(b)}=D\otimes\mathbf1\mathbf1^t/b$ (K08:2--3), y $D^{(b)}-J_{bk}=B\otimes\Pi_b$ con
  $\Pi_b=\mathbf1\mathbf1^t/b$ proyección ortogonal. Así $bk\max D^{(b)}=k\max D$ y
  $\|B\otimes\Pi_b\|_{\rm op}=\|B\|_{\rm op}$: la clase no cambia y $K=bk$ no está acotado. Aquí la
  uniformidad en $K$ es esencial: el límite $b\to\infty$ de K08:8--9 necesita una constante
  independiente de $b$, y el Teorema la da.
  \item E01, a través de $R=g_T-u_T$ (K06:5), en los menores $A^-_c$ y en la extensión con cruce
  ausente y hasta cuatro supresiones (K08:11--52; W2t:135--166). La clase se hereda del padre
  por D04:52--63 (orig.~45--56), K02 y U02:35--45 (E1; W2t:148--159, 225); $K\in[k-4,k]$, y sus
  replicaciones.
  \end{itemize}
\item \emph{M01} (M01:18--23; MM07:2--14): $u_T+R_T$ en los menores $D_{-X}$, con $|X|\le4$, del
  exterior. Es el caso anterior (W2t:225).
\item \emph{GI02}:61--66: variación de $u_T+R_T$ entre $D_0$ y $(D_0)_{-X}$, con $|X|\le v$. La clase
  la dan GI02:4--5 y U02:35--45 (E1): un $r'<1$ común si $v/j$ es pequeño. Vale también con
  $v\le8\log n$ (WG:56), porque $C_*$ no depende de $v$.
\end{enumerate}
El Teorema sólo se usa en datos físicos reales. Las prolongaciones complejas no lo usan: $\phi$ en
$\mu_h$ (MM07:16--44), $\Gamma_z$ en discos (K08:19--20) y $H_*$ usan $H_2$ y U02--U04. El resto
$O(s^3)$ de \path{07_coeficiente.tex}:54--55 de la ronda~9 sale de U04:44--46 con $q=2$
(W2t:106), no del Teorema.

\paragraph{3. Dependencias.} El Teorema usa sólo el nodo N1 (U01--U05; WG:168) y herramientas
estándar. La arista DW$\to$N6 de WG:191 se sustituye por N1$\to$N6, que WG:192 ya registra como
«ruta reparadora». El grafo sigue acíclico, con el mismo orden topológico (WG:205).

\paragraph{4. Qué queda de Divoux--Wang.} Nada en esta prueba. U01 con $h=v=0$ usa cuatro
herramientas:
\begin{itemize}
\item la identidad de mezcla (01m:56--58; E1:86);
\item la fórmula de Wick (01m:35--42; E1:75);
\item la cota $|R_B|\le e^{C\sqrt K}$ en el disco, con Cauchy (01m:23--32, 51--53; E1:74, 77);
\item la cola gamma ponderada (01m:59--66; E1:87--91).
\end{itemize}
Son los cuatro elementos de U01:37--39, rehechos por E1 (E1:232; WG:142). No se usan ni el Lema~3.7,
ni las cotas (5)--(8), ni la suma binomial, ni la Proposición~3.1 de Divoux--Wang (Wp:32--66).
